\chapter{Online Map Refinement}
\label{ch:refiner}

\section{Design}
\label{sec:method:refiner}

Prediction supplies an initialisation. A separate process photometrically
refines the assembled map \emph{while SLAM runs}; together with the GUI,
the complete system has four processes (\cref{fig:overview}).

\subsection{Parameterisation}
Gaussians are optimised in their anchor keyframe's local frame in
pre-activation form --- $\log s$, unnormalised $q$, $\mathrm{logit}\,\alpha$,
and the SH DC term --- so that scales stay positive and opacity stays in
$(0,1)$ without constraints, and composed to the world by
\cref{eq:anchor} at every step. Per-group Adam rates follow the 3DGS
reference, with the positional rate scaled by scene extent:
$\eta_\mu = 1.6\!\times\!10^{-4}\!\cdot\!\mathrm{extent}$,
$\eta_{f_{dc}} = 2.5\!\times\!10^{-3}$, $\eta_\alpha = 5\!\times\!10^{-2}$,
$\eta_s = 5\!\times\!10^{-3}$, $\eta_q = 1\!\times\!10^{-3}$.

\subsection{Objective}
Supervision views are tracked frames, so the refiner optimises
\begin{equation}
\mathcal{L}_{\text{ref}} = (1-\lambda_{D})\,\lVert \hat{I}-I \rVert_1
                          + \lambda_{D}\,\bigl(1-\mathrm{SSIM}(\hat{I},I)\bigr),
\qquad \lambda_{D}=0.2 ,
\label{eq:refloss}
\end{equation}
i.e.\ the 3DGS objective rather than \cref{eq:splatt3rloss}: LPIPS is a
reporting metric here, not a training signal, so that our evaluation metric is
not also our loss.

\subsection{Anchor-carried supervision}
Each supervision frame is stored as an image together with
$(\text{anchor index}, T_{\text{anchor},f})$ rather than a world pose. Its
world pose is recomposed through the anchor's \emph{current} pose at sample
time, so supervision follows loop closures exactly as the map does under
\cref{eq:anchor,eq:supervisionpose}. Frames are held in CPU shared memory,
which allows the refiner to run on a second GPU.
The refiner consequently never becomes a second pose-graph writer: a loop
closure changes the anchor pose once, and both map placement and supervision
views follow that correction.

\subsection{Sampling}
Supervision is drawn from a reservoir of $200$ frames spanning the whole
sequence plus a ring of the $64$ most recent, with a recent fraction of $0.3$.
We measured uniform-over-history to dominate a recency-weighted schedule
(held-out early/late split: $14.40/14.09/14.72$ uniform versus
$14.06/13.79/14.35$ recent-only), because causal arrival already biases the
pool toward recent views.

\subsection{Scheduling}
Sharing one GPU, the refiner holds a duty cycle $d$ by sleeping
$\mathrm{EMA}[\Delta t]\,(1-d)/d$ after each step, so optimisation steps are
dropped but never frames. Refiner load on the \emph{same} card doubles tracker
latency ($101\!\rightarrow\!206$\,ms p50); on a second card it is noise
($103$\,ms), which is the configuration we deploy. After the sequence ends an
optional polish phase continues with the final poses.

\subsection{Centre updates}
The external-comparison configuration holds $\mu$ fixed. Earlier fixed-map
experiments also favoured this choice at approximately $300$ and $3000$
steps. The broader causal delivery study did not show the same preference.
External comparisons and causal studies therefore report their centre
settings separately from the command-line defaults in \cref{app:hyper}.

When centres are fixed, the positional learning rate above is inactive.
When enabled, centre updates can trade geometric consistency against fit
to supervision cameras. \Cref{sec:exp:online} separates the tested causal
configuration from the fixed-centre, post-sequence comparison.

The refiner publishes a versioned CPU snapshot rather than mutable GPU state.
The visualiser uploads only complete versions and retains the previous valid
version on a read/write collision. If the display buffer is smaller than the
refined map, publication subsamples after the full map has been saved. This
separates bounded display memory from the evaluation artifact and prevents a
large-map publication failure from discarding a valid refinement result.

\section{The polish phase}
\label{sec:refiner:polish}

When the sequence ends, the trajectory stops changing but the map is still
under-optimised: every step taken during the sequence was taken against poses
that were still moving. A \emph{polish} phase continues refinement with the
final poses, and it is the single largest online win we measured:
$+2.21$\dB{} PSNR and $-16.7\%$ LPIPS over seven TUM \texttt{fr1} sequences.

The required duration varies. Sequences need
between $1445$ and $3821$ steps to converge --- a $2.6\times$ spread --- and no
static quantity we tried predicts which sequence needs which: not length, not
keyframe count, not final Gaussian count, not the loss at the start of the
phase. A fixed step budget therefore either wastes time on the easy sequences
or truncates the hard ones. The implementation offers a convergence
criterion (\texttt{--refiner-polish-tol} with a patience of two logging
windows), disabled by default. Its stopping rule and fixed-duration
experiments should be reported separately.

\Cref{tab:tum} uses a fixed $300$\,s budget for our external comparison,
rather than convergence-based stopping. Its PSNR is close to MonoGS after
$339$\,s colour refinement, while MonoGS has lower LPIPS. A separate
historical $1200$\,s run reached $14.61$\,dB and $0.379$ LPIPS under the
older evaluation configuration. That experiment demonstrates further
optimisation headroom; it is not mixed into the fresh full-SH comparison or
presented as a gain at matched total compute.

\section{Density control}
\label{sec:refiner:omissions}

The refiner reinstates the 3DGS objective (\cref{eq:3dgsloss}) but not the
adaptive density control that accompanies it in the reference implementation.
This choice has a measurable map-size cost.

Cloning and splitting are omitted because the map already has an excess of
primitives rather than a shortage: the failure mode this thesis addresses is
crowding, and a densification policy tuned for a sparse SfM initialisation
would be solving the opposite problem. Pruning and periodic opacity reset
are also omitted. Pruning changes the primitive count and requires
corresponding updates to optimiser and snapshot state. Opacity reset changes
the optimisation schedule.

The cost is visible in \cref{sec:lim:compact}: without pruning, nothing ever
removes a primitive that the data does not support, and the map size is
whatever the injection pipeline produced. A voxel-based deduplication pass was
implemented as a partial substitute and measured as a negative: it removes
$27.6\%$ of the map for $-0.53$\dB{}. We report it as a map-size control that
costs quality when it fires without a recovery budget, and it is off by
default. Refining directly under a primitive budget --- rather than pruning
after the fact --- is the future work named in \cref{ch:conclusion}.

\section{Tracking preservation}
\label{sec:refiner:invariant}

The refiner reads and composes poses through
\cref{eq:anchor,eq:supervisionpose}, without optimising them. The recorded
ablations preserve ATE ($0.016975$\,m, rounding to the $0.0170$\,m baseline)
while changing map quality by several dB. This supports the intended
separation of mapping and tracking. It does not eliminate resource contention
or guarantee identical scheduling in every multi-process run.
